\scanpage{45}
% archive sheet number: 44
\mspage{78}
\hfill n\textsuperscript{o} 382

\[
P^k_x:C^r_x(X)\longrightarrow P^k_x(X)
\]
soit un homomorphisme ; on a les formules :
\[
P^k_x(f)+P^k_x(g)=P^k_x(f+g)\;,\quad\lambda.P^k_x(f)=P^k_x(\lambda.f)\;,\quad P^k_x(f).P^k_x(g)=P^k_x(f.g)\;,
\]
pour $f,g\in C^r_x(X)$, $\lambda\in K$.

7.1.3. Soit de plus $h$ un morphisme dans une variété $Z$ de classe $C^r$ d'un voisinage de $y$. Alors $j^k_x(h\circ f)$ ne dépend que de $j^k_x(f)$ et de $j^k_y(h)$ et on l'appelle le \uline{produit} (ou le \uline{composé}) de $j^k_x(f)$ et $j^k_y(h)$ ; le composé $v\circ u$ de deux jets $u$ et $v$ est donc défini si le but de $u$ coïncide avec la source de $v$ ; ce produit est associatif. Un jet $u$ d'ordre $k$ de $X$ dans $Y$ est \uline{inversible} s'il existe un jet $v\in J^k(Y,X)$ tel que $v\circ u$ et $u\circ v$ soient définis et égaux aux jets des applications identiques $1_X$ et $1_Y$. Si $k=0$, tout jet est inversible ; si $k\geq 1$, un jet $u$ est inversible si et seulement si c'est la partie principale d'ordre $k$ d'un isomorphisme local de $X$ dans $Y$ ; dans ce cas, tout morphisme $f$ tel que $j^k_x(f)=u$ est un isomorphisme local en $x$. Un jet d'ordre $k$ inversible de $X$ dans $K^p$ (resp. de source $x$) est appelé un \uline{repère d'ordre} $k$ de $X$ (resp. de source $x$). L'ensemble de ces repères est noté $R^k(X)$ (resp. $R^k_x(X)$). Si $U$ et $V$ sont des ouverts de $X$ et $Y$, on a des inclusions canoniques
\[
J^k(U,V)\subset J^k(X,Y)\qquad\text{et}\qquad R^k(U)\subset R^k(X)\;.
\]

7.1.4. Si $k'\leq k$, le jet $j^{k'}_x(f)$ d'ordre $k'$ d'un germe de morphisme $f$ ne dépend que de $j^k_x(f)$ et est inversible si $j^k_x(f)$ l'est. En faisant correspondre $j^{k'}_x(f)$ à $j^k_x(f)$, on définit des applications canoniques surjectives
\[
J^k_x(X,Y)_y\longrightarrow J^{k'}_x(X,Y)_y\;,\qquad R^k_x(X)\longrightarrow R^{k'}_x(X)
\]
et $P^k_x(X)\longrightarrow P^{k'}_x(X)$ ; la dernière est un homomorphisme d'algèbres.

7.1.5. Soit $f=(f^1,\ldots,f^q)$ une application de classe $C^r$ d'un ouvert $U$ de $K^p$ dans l'espace $K^q$. Si $x\in U$, le jet $j^k_x(f)$ de $f$ est complètement déterminé par les coordonnées de sa source $x$ et par les valeurs en $x$ des fonctions $\Delta^\alpha f^i$, avec
